\section*{学习导航：向量、凸性与无穷远方向}
本讲对应 Debreu《Theory of Value》第一章第9节。课件后半部分补充了渐近锥理论，并重复前半部分若干证明。原页全部保留；重复页仍逐页说明其作用，并指向首次完整论证。先掌握第3--23页的运算，再学第46--66页的凸性，最后回到渐近锥，是较适合初学者的阅读顺序。

\paragraph{先认符号。} $\R^N$ 是所有 $N$ 个实数组成的有序数组；$x=(x_1,\ldots,x_N)$，$\norm{x}=(\sum_i x_i^2)^{1/2}$ 是到原点的距离。$x_k\to x$ 表示距离 $\norm{x_k-x}\to0$。下标 $k$ 常是序列编号，分量编号另用 $i$，不要混为一谈。$\overline C$ 表示闭包，即 $C$ 内点列所有可能的极限组成的集合；$\operatorname{int}C$ 表示能在 $C$ 中放下一个小开球的点。

\paragraph{本讲反复使用的四件工具。} 三角不等式 $\norm{u+v}\le\norm u+\norm v$；有限维有界序列有收敛子列；闭集包含其内部点列的全部极限；在 $\R^N$ 中紧集恰为闭且有界的集合。注意“有界”只保证某个子列收敛，绝不保证原序列收敛。写证明时必须保留“取子列”这一步。

\paragraph{渐近锥的约定。} 本文将课件 $AC$、$C_\infty$ 统一理解为非空集合的渐近锥，且总含 $0$；$0^+C$ 专指衰退锥。它们对非空\emph{闭凸集}相等，对一般集合不相等。遇到空集合的交、和时先单独处理，避免套用要求非空的序列定义。

\Slide{1}{本讲研究什么}
本讲标题是“$\R^N$ 中的向量”，但内容远超向量计算：研究集合相加是否保持闭性、凸性，以及无界集合往哪些方向延伸。经济学中多个生产集合相加就是总生产可能集合；证明这个总集合闭，才能保证极限方案仍可行。封面给出教师、教材与授课时间，不含需要证明的命题。

\Slide{2}{在课程知识链中的位置}
前面的极限、连续性和对应理论，本讲都要用到。主线是：向量加法连续 $\Rightarrow$ 紧集之和紧；无穷远方向不能互相抵消 $\Rightarrow$ 闭集之和闭；凸性 $\Rightarrow$ 可用超平面描述边界。下一讲的不动点定理又需要紧性与凸性。本页的目录不是一组定理，先用它定位依赖关系。

\Slide{3}{向量相加：逐个坐标相加}
若 $x^1=(1,2)$、$x^2=(3,-1)$，则 $x^1+x^2=(4,1)$，$-x^1=(-1,-2)$，$x^1-x^2=(-2,3)$。上标 $1,2$ 标记两个向量，并非幂。一般写 $\sum_{k=1}^K x^k$，其第 $i$ 个坐标是 $\sum_{k=1}^K x_i^k$。课件末行的求和指标排版有混用，记住“$k$ 数向量、$i$ 数坐标”即可。零向量的每个坐标都是 $0$。

\Slide{4}{数乘及平行四边形直观}
$t x=(t x_1,\ldots,t x_N)$：$t>1$ 拉长，$0<t<1$ 缩短，$t<0$ 还反向，$t=0$ 变为原点。向量加法可看成先走 $x^1$ 再走 $x^2$；终点是平行四边形的对角顶点。若 $x=0$，所谓“直线 $0x$”不唯一，但代数定义仍给 $t0=0$，不受影响。

\Slide{5}{为什么加法和数乘连续}
\paragraph{加法的完整验证。} 若 $u_k\to u$、$v_k\to v$，则
\[\norm{(u_k+v_k)-(u+v)}\le\norm{u_k-u}+\norm{v_k-v}\longrightarrow0.\]
\paragraph{数乘的完整验证。} 若 $t_k\to t$、$x_k\to x$，把差拆成 $t_k(x_k-x)+(t_k-t)x$。因为收敛数列 $(t_k)$ 有界，存在 $M$ 使 $|t_k|\le M$，于是
\[\norm{t_kx_k-tx}\le M\norm{x_k-x}+|t_k-t|\norm x\longrightarrow0.\]
这是课件“显然连续”背后的理由。考试不能只说“极限可以交换”，应给出连续性或以上估计。

\Slide{6}{点与向量是两种看法}
同一个数组 $(x_1,\ldots,x_N)$，说“点”时强调位置，说“向量”时强调从原点到该点的位移。比如 $x-y$ 是从 $y$ 指向 $x$ 的位移。后面的集合平移和法向量都用这种转换；不需要另学一种新的对象。

\Slide{7}{平移}
固定 $a$，$T_a(x)=x+a$ 将所有点沿同一方向搬动。它连续，逆映射 $T_{-a}$ 也连续，因此保留开集、闭集、紧集、连通性。证明闭性也可直接做：若 $x_k+a\to z$ 且 $x_k\in C$，则 $x_k\to z-a$；$C$ 闭推出 $z-a\in C$，故 $z\in C+a$。平移不改变无穷远方向，见第31页。

\Slide{8}{集合相加：闵可夫斯基和}
$A+B=\{a+b:a\in A,b\in B\}$，意为两边各挑一个点，再把所有可能的和收集起来。它与并集 $A\cup B$ 完全不同。例如 $[0,1]+[2,4]=[2,5]$，而并集是两段。证明 $z\in A+B$ 时，应展示分解 $z=a+b$；证明不存在这种分解即可证明不属于和集。

\Slide{9}{两条线段相加得到正方形}
课件中的第一条线段是 $A=\{(s,0):0\le s\le1\}$，第二条是 $B=\{(0,t):0\le t\le1\}$。任取两点相加得到 $(s,t)\in[0,1]^2$，故 $A+B\subset[0,1]^2$；反过来每个 $(s,t)$ 都可写为 $(s,0)+(0,t)$，故另一包含也成立。集合等式通常要这样证明两个方向，不能只凭图形判断。

\Slide{10}{集合取负、作差与闭包的和}
$-A=\{-a:a\in A\}$，$A-B=A+(-B)$，并非集合差 $A\setminus B$。本页命题正确读作
\[\overline{A_1}+\cdots+\overline{A_K}\subseteq\overline{A_1+\cdots+A_K}.\]
任取左边的 $z=\sum_k a_k$。因为 $a_k\in\overline{A_k}$，对每个 $k$ 选 $a_{k,q}\in A_k$ 使 $a_{k,q}\to a_k$。有限个极限可以相加，故 $z_q=\sum_k a_{k,q}\to z$，且 $z_q$ 在和集内。按闭包定义，$z$ 在右边。横线所在位置至关重要；这并没有说闭集之和一定闭。

\Slide{11}{紧集之和紧：先看有界性}
若 $A_k$ 紧，则有常数 $M_k$ 使 $\norm{a_k}\le M_k$。任取 $z\in\sum_k A_k$，有 $z=\sum_k a_k$，从而 $\norm z\le\sum_k\norm{a_k}\le\sum_k M_k$。上界不依赖所选 $z$，故整个和集有界。如果有一个 $A_k=\varnothing$，和集为空，紧性自动成立；下面默认各集合非空。此页只完成“有界”，还不能宣告紧。

\Slide{12}{紧集之和紧：连续像的短证明}
有限乘积 $P=A_1\times\cdots\times A_K$ 紧，映射 $F(a_1,\ldots,a_K)=\sum_k a_k$ 连续，而 $F(P)=\sum_k A_k$。连续映射把紧集送到紧集，所以结论成立。这条论证其实已同时得到闭性和有界性。
\paragraph{为什么不能把紧换成闭。} 在 $\R^2$ 中取 $A=\{(u,1/u):u>0\}$、$B=\{(v,0):v\in\R\}$，两者都闭，然而 $A+B=\R\times(0,\infty)$ 不闭。$A$ 的闭性来自：若它的点列有有限极限，则第一坐标不可能趋于 $0$ 或无穷，极限仍满足乘积为 $1$ 且正。

\Slide{13}{紧集之和紧：逐层抽取子列}
不用连续像定理也能证明。给定 $z_q=\sum_{k=1}^K a_{k,q}$，先利用 $A_1$ 紧，挑索引子列，使第一个分量收敛；\emph{只在这条子列中}再挑，使第二个分量收敛。第一个分量仍然收敛，因为收敛序列的子列保留原极限。重复有限次，得到同一条索引子列 $q_j$，使每个 $a_{k,q_j}\to a_k\in A_k$。不能分别随意选 $K$ 条不同索引子列后直接相加。

\Slide{14}{紧集之和紧：把子列重新相加}
沿第13页得到的共同索引，
\[z_{q_j}=\sum_{k=1}^K a_{k,q_j}\longrightarrow\sum_{k=1}^K a_k\in\sum_{k=1}^K A_k.\]
于是和集的每个序列都有极限仍在和集内的收敛子列，正是有限维空间中的序列紧性。这页收尾的逻辑顺序是“任取序列—构造子列—验证极限留在集合”，不是“找到了某个收敛序列所以紧”。

\Slide{15}{有限个连续函数相加仍连续}
设 $f=\sum_{k=1}^K f_k$，要证在 $x_0$ 连续。给定 $\eps>0$，每个 $f_k$ 连续，选 $\delta_k>0$ 使 $\norm{x-x_0}<\delta_k$ 时 $\norm{f_k(x)-f_k(x_0)}<\eps/K$。取 $\delta=\min_k\delta_k>0$，三角不等式给
\[\norm{f(x)-f(x_0)}\le\sum_k\norm{f_k(x)-f_k(x_0)}<\eps.\]
“有限”保证最小值仍为正；无限级数不能直接照抄，通常要增加一致收敛等条件。

\Slide{16}{对应相加：为什么需要控制值域}
对应 $\varphi_k(x)$ 给出一个可选集合；$\varphi(x)=\sum_k\varphi_k(x)$ 给出所有总选择。课件把各值域放在紧集 $T_k$ 中，以便从分解的各分量抽取收敛子列。采用本教材的紧值、闭图式上半连续约定时，上半连续性得以保持。若使用一般拓扑教材中不自动要求闭值的“开集定义”，须另外明确紧值等条件，不能不加条件地把闭图与上半连续等同。下半连续表示旧选择能由邻近参数下的选择逼近；求和的下半连续不需要各 $T_k$ 紧。

\Slide{17}{对应之和半连续：补全漏掉的抽子列步骤}
\paragraph{上半连续。} 设 $x_q\to x_0$，$y_q\in\varphi(x_q)$，$y_q\to y$。分解 $y_q=\sum_k y_{k,q}$，其中 $y_{k,q}\in\varphi_k(x_q)\subset T_k$。由各 $T_k$ 紧，\emph{取共同子列}使 $y_{k,q_j}\to y_k$。各对应在 $x_0$ 的闭图性质给 $y_k\in\varphi_k(x_0)$，故 $y=\sum_k y_k\in\varphi(x_0)$。总值域紧，因此这推出所需上半连续性。原页直接写各分量收敛，省略了关键依据。
\paragraph{下半连续。} 给定 $y=\sum_k y_k\in\varphi(x_0)$ 和任意 $x_q\to x_0$。各 $\varphi_k$ 下半连续，存在 $y_{k,q}\in\varphi_k(x_q)$ 且 $y_{k,q}\to y_k$；取和便得 $y_q\in\varphi(x_q)$ 且 $y_q\to y$。这里要逼近的是\emph{指定的}旧选择 $y$。

\Slide{18}{加权平均与凸组合的区别}
$t_1+t_2=1$ 时 $t_1x^1+t_2x^2$ 称仿射组合，课件称 weighted average。若再有 $t_1,t_2\ge0$，才是凸组合。比如 $2x^1-x^2$ 的系数和为 $1$，却可能跑到线段外。后面证明凸性时必须检查“系数非负”和“系数和为 $1$”两个条件。

\Slide{19}{直线、射线：参数的范围决定形状}
统一公式是 $z=x+t(y-x)$。$t\in\R$ 给整条直线；$t\ge0$ 给从 $x$ 出发经过 $y$ 的闭射线；$t>0$ 去掉起点，给开射线。取 $x=(1,1)$、$y=(3,2)$，则 $z=(1+2t,1+t)$。当 $t=0,1,2$ 时分别是起点、$y$、继续向前的点。“闭射线”指包含端点，并不是有界。

\Slide{20}{图5：方向相同，起点不同}
图中下方射线从 $0$ 出发，方向是 $y-x$；上方射线由下方整条平移 $x$ 得到，起点变成 $x$ 并经过 $y$。因此“方向”不包含起点信息。图里的斜率在二维画面可看见，但 $\R^N$ 中更可靠的描述始终是 $x+t(y-x)$。

\Slide{21}{线段与退化线段}
把第19页的参数限制为 $0\le t\le1$，得到 $[x^1,x^2]$。两端都包括在内。若 $x^1=x^2$，任意 $t$ 都给同一个点，线段退化成单点。以后用“包含任意两点的整段”定义凸性时，单点也应允许，不能强行要求两点不同。

\Slide{22}{锥：沿射线可以一直走}
顶点为 $x$ 的锥 $C$ 要满足：$x\in C$，且任意 $y\in C$、$t\ge0$ 都有 $x+t(y-x)\in C$。顶点为 $0$ 时简化为 $t y\in C$。\emph{锥不一定凸}：二维正横轴与正纵轴的并是锥，却不包含 $(1,0)$ 与 $(0,1)$ 的中点。凸锥还要容许不同方向的非负线性组合。

\Slide{23}{正半独立：方向不能抵消}
锥 $C_1,\ldots,C_K$ 正半独立，意思是 $x_k\in C_k$ 且 $\sum_kx_k=0$ 只能发生在全部 $x_k=0$ 时。两个锥时等价于 $C_1\cap(-C_2)=\{0\}$。
\paragraph{两方向证明。} 若交仅含 $0$，由 $x_1+x_2=0$ 得 $x_1=-x_2\in C_1\cap(-C_2)$，所以两者均零。反过来，若正半独立，任取 $z\in C_1\cap(-C_2)$，则 $z\in C_1$、$-z\in C_2$，两者和为零，故 $z=0$。两个同向非负射线正半独立，两个相反射线则不满足。

\Slide{24}{渐近锥：只保留越来越远的点}
先截取 $S^M=\{x\in S:\norm x\ge M\}$，再作包含它的最小闭锥 $\Gamma(S^M)$，最后取所有 $M\ge0$ 的交。这样的方向必须在无论多远的截尾中都能被逼近。锥在这里总含原点，特别约定 $\Gamma(\varnothing)=\{0\}$。每个 $\Gamma(S^M)$ 闭且对非负数乘封闭，交仍闭且是锥，所以 $AS$ 是闭锥。它不一定凸，除非再对 $S$ 加凸性条件。

\Slide{25}{为什么截尾锥递减；闭包不能漏}
若 $M_1\le M_2$，距离门槛提高，故 $S^{M_2}\subset S^{M_1}$，进而 $\Gamma(S^{M_2})\subset\Gamma(S^{M_1})$。所谓极限就是递减集合族的交。准确的公式是
\[\Gamma(E)=\overline{\{\alpha x:\alpha\ge0,\ x\in E\}\cup\{0\}}.\]
闭包负责加入仅能逼近的边界方向；$\{0\}$ 负责空集情形。去掉闭包、或把“存在任意接近的点”误写成“恰好存在同方向的点”，都会使后面的证明失效。

\Slide{26}{序列定义：把远点缩小后取极限}
对非空 $C$，$d\in C_\infty$ 等价于存在 $x_k\in C$、$t_k\to+\infty$ 使 $x_k/t_k\to d$。这里可以删掉有限个初项，使 $t_k>0$。例：$C=\{(s,1):s\ge0\}$ 的渐近锥为 $\{(a,0):a\ge0\}$，因为缩放后第二坐标 $1/t_k\to0$。零方向总存在：固定一点 $x\in C$，用 $x/k\to0$。非零方向才表示真正逃向无穷远。

\Slide{27}{衰退锥：从每个点都能沿它走}
定义 $0^+C=\{d:x+\lambda d\in C\text{ 对所有 }x\in C,\lambda\ge0\}$。这比“存在远点列沿该方向”强得多。一般只有 $0^+C\subset C_\infty$；$C$ 非空闭凸时才可保证相等，完整证明见第110--111页。\textbf{条件提示：} 本页末行 $AC=0^+C$ 若只假设凸而未假设闭，并不普遍成立；第112页正好给反例。不要把不同作者的术语相同误当作无条件等价。

\Slide{28}{渐近锥性质总览与两个边界条件}
正确读法：$C_\infty$ 闭且是锥；$(\overline C)_\infty=C_\infty$；若 $C$ 本身是锥，则 $C_\infty=\overline C$，横线不能漏。证明在第77--80页。
归一化方向集 $C_N$ 是无界点列 $x_k$ 的单位方向 $x_k/\norm{x_k}$ 的极限集。为兼容有界 $C$ 导致 $C_N=\varnothing$ 的情形，等式应写
\[C_\infty=\{0\}\cup\{\lambda u:\lambda>0,\ u\in C_N\}.\]
“$C$ 有界当且仅当 $C_\infty=\{0\}$”按本页上下文要求 $C$ 非空，空集另行处理。

\Slide{29}{有界当且仅当没有非零渐近方向}
若 $\norm x\le M$ 对所有 $x\in C$ 成立，那么任意 $x_k/t_k$ 的范数至多 $M/t_k\to0$，故只可能得到 $0$；非空性保证 $0$ 确实存在。
反向用逆否命题：若 $C$ 无界，对每个整数 $k$ 选 $x_k\in C$ 使 $\norm{x_k}\ge k$。单位向量 $u_k=x_k/\norm{x_k}$ 位于紧的单位球面，所以有子列趋向 $u$，且 $\norm u=1$。取 $t_j=\norm{x_{k_j}}\to\infty$，就得到 $u\in C_\infty$，是非零方向。关键是“无界—归一化—球面紧—非零极限”。

\Slide{30}{单调性、平移不变性与乘积的预告}
$S_1\subset S_2$ 时，在 $S_1$ 中见证某渐近方向的点列也在 $S_2$ 中，所以 $AS_1\subset AS_2$。若 $T\ne\varnothing$，取固定 $a\in T$，则 $S+a\subset S+T$；结合平移不变性，$AS=A(S+a)\subset A(S+T)$。这里 $a$ 必须固定，不能随序列任意变化。乘积的结论一般只是包含，原因是各分量需要共同的缩放尺度，见第33页。

\Slide{31}{平移不改变渐近锥：逐步证明}
设 $d\in(S+a)_\infty$，存在 $z_k\in S+a$、$t_k\to\infty$ 使 $z_k/t_k\to d$。写 $z_k=x_k+a$，$x_k\in S$，则
\[\norm{x_k/t_k-d}\le\norm{z_k/t_k-d}+\norm a/t_k\to0.\]
故 $(S+a)_\infty\subset S_\infty$。把 $S+a$ 当作新集合，再平移 $-a$，得到反向包含。直观是固定距离的搬动除以越来越大的尺度后消失。

\Slide{32}{平移不变性的等价链写法}
本页将第31页的双包含压缩为等价关系。核心等式是 $(x_k+a)/t_k=x_k/t_k+a/t_k$，最后一项趋于 $0$，所以两边具有同一个极限。练习时先自行写出 $x_k$ 属于哪个集合，再写极限；不必背每一行排版。若平移量变为 $a_k$ 且 $a_k/t_k$ 不趋零，结论便可能改变。

\Slide{33}{乘积的渐近锥包含于各渐近锥的乘积}
设 $C=S_1\times\cdots\times S_K$ 非空，$d\in C_\infty$。取 $x_q=(x_{1,q},\ldots,x_{K,q})\in C$、$t_q\to\infty$ 使 $x_q/t_q\to d=(d_1,\ldots,d_K)$。向量收敛蕴含每个坐标块收敛，所以 $x_{k,q}/t_q\to d_k$，即 $d_k\in(S_k)_\infty$。因此
\[C_\infty\subset\prod_{k=1}^K(S_k)_\infty.\]
反向不自动成立：各个方向的见证序列可能只在互不匹配的尺度出现。对凸集，可由第105--108页的渐近正则性解决尺度匹配问题。

\Slide{34}{两个核心应用：交有界，和闭}
交集有界的充分条件是所有渐近锥的交仅为 $\{0\}$；闭集之和闭的充分条件是它们的渐近锥正半独立。两者都在阻止某种“逃到无穷”的行为：前者阻止同时在所有集合中逃走，后者阻止各分量逃走但总和相互抵消。它们是\emph{充分条件}，不要写成必要条件。例如整个 $\R$ 加自己仍闭，两个渐近锥却能抵消。

\Slide{35}{交集有界：用单调性证明}
令 $D=\bigcap_{\gamma\in I}C_\gamma$。若 $D=\varnothing$，有界性直接成立。若非空，则 $D\subset C_\gamma$ 给 $D_\infty\subset(C_\gamma)_\infty$ 对每个 $\gamma$ 成立，因此
\[D_\infty\subset\bigcap_\gamma(C_\gamma)_\infty=\{0\}.\]
因为 $D$ 非空，$0\in D_\infty$，于是 $D_\infty=\{0\}$，第29页即给 $D$ 有界。这就是“先处理空集，再使用非空定理”的标准写法。

\Slide{36}{不能从递减交直接跳到某个截尾为空}
本页用截尾集合重述第35页。结论正确，但中间“存在 $M_0$ 使截尾为空”需要第29页的有界性判别，不能只由递减集合交为空或交为 $\{0\}$ 直接断言。一般递减非空集可以有空交，例如 $[n,\infty)$。在这里，可反证：若交集无界，就抽出单位极限方向，同时属于全部渐近锥，与仅含零矛盾。

\Slide{37}{线性映射何时把闭集送到闭集}
线性映射 $L$ 满足 $L(\alpha x+\beta y)=\alpha L(x)+\beta L(y)$；有限维时它连续。核 $\ker L=L^{-1}(\{0\})$ 是所有被压到零的方向。条件
\[\ker L\cap C_\infty=\{0\}\]
表示 $C$ 的非零逃逸方向不能被 $L$ 完全压扁。若 $C$ 闭，则它保证 $L(C)$ 闭，而且 $L(C_\infty)=(L(C))_\infty$。无条件成立的方向只有 $L(C_\infty)\subset(L(C))_\infty$。下面四页分别处理“防逃逸、闭性、正向包含、反向包含”。

\Slide{38}{线性像证明的核心：归一化反证}
若存在 $x_k\in C$，$\norm{x_k}\to\infty$，而 $L(x_k)$ 有界，则单位向量取子列 $x_k/\norm{x_k}\to u$，$\norm u=1$。按定义 $u\in C_\infty$；连续性和线性给
\[L(u)=\lim_k \frac{L(x_k)}{\norm{x_k}}=0.\]
于是 $u\in\ker L\cap C_\infty$ 且非零，与条件矛盾。更准确的可用结论是：$C$ 中任何像有界的点列自身必有界；若自身无界，先抽取范数趋无穷的子列，再运用本证明。

\Slide{39}{线性像闭：把收敛提升回原集合}
任取 $y_k\in L(C)$ 且 $y_k\to y$，每次选一个原像 $x_k\in C$ 满足 $L(x_k)=y_k$。$(y_k)$ 有界，第38页使 $(x_k)$ 有界；抽取 $x_{k_j}\to x$。$C$ 闭给 $x\in C$，连续性给 $L(x)=\lim_j y_{k_j}=y$，故 $y\in L(C)$。注意 $L$ 可能不单射，原像无须唯一；任意选取的原像都受这里的有界性论证控制。

\Slide{40}{渐近方向先取像：总是可行}
任取 $d\in C_\infty$，选 $x_k\in C$、$t_k\to\infty$ 使 $x_k/t_k\to d$。由于 $L(x_k)\in L(C)$ 且
\[\frac{L(x_k)}{t_k}=L\!\left(\frac{x_k}{t_k}\right)\longrightarrow L(d),\]
得 $L(d)\in(L(C))_\infty$。这一步既不使用 $C$ 闭，也不使用核条件。考试中指出哪些假设用在哪里，能避免把强条件当成所有步骤的理由。

\Slide{41}{线性像渐近锥的反向包含}
须补严谨步骤，不能直接把第38页用于不在 $C$ 内的 $x_n/t_n$。若 $Lx_n/t_n\to y$，假设 $x_n/t_n$ 无界，则抽子列使 $r_n=\norm{x_n}/t_n\to\infty$。归一化后再抽 $x_n/\norm{x_n}\to d$，其中 $\norm d=1$，$d\in C_\infty$ 且 $Ld=\lim (Lx_n/t_n)/r_n=0$，违反核条件。因此 $x_n/t_n$ 有界，抽子列趋于 $d$，定义给 $d\in C_\infty$ 且 $Ld=y$。

\Slide{42}{闭集和与正半独立}
正半独立指若 $d_i\in(C_i)_\infty$ 且 $\sum_i d_i=0$，则每个 $d_i=0$。设积集 $C=\prod_iC_i$，线性映射 $L(x_1,\ldots,x_m)=\sum_i x_i$。逐坐标极限给 $C_\infty\subseteq\prod_i(C_i)_\infty$，所以核条件成立。前面的线性像定理给 $L(C)=\sum_iC_i$ 闭，且 $(\sum_i C_i)_\infty=L(C_\infty)\subseteq\sum_i(C_i)_\infty$。这说明渐近方向不能互相抵消是保证闭和的重要条件；闭集的和一般不自动闭。

\Slide{43}{闭集和与正半独立}
本页继续“闭集和与正半独立”。完整定义、条件和论证见本讲原页第 42 页下方。

\Slide{44}{闭集和与正半独立}
本页继续“闭集和与正半独立”。完整定义、条件和论证见本讲原页第 42 页下方。

\Slide{45}{闭集和与正半独立}
本页继续“闭集和与正半独立”。完整定义、条件和论证见本讲原页第 42 页下方。

\Slide{46}{凸集与多面锥}
凸指任取 $x,y\in C$ 及 $0\le t\le1$，$tx+(1-t)y\in C$。这是整条线段都在集合中；不是只检验中点。多面锥由有限向量的非负线性组合构成，权重不要求和为1；凸组合则要求和为1。二者必须区分。

\Slide{47}{凸集与多面锥}
本页继续“凸集与多面锥”。完整定义、条件和论证见本讲原页第 46 页下方。

\Slide{48}{凸性的基本运算及连通性}
交：两端点属于每个凸因子，线段逐个属于它们，故属于交。和：写 $x=\sum x_i,y=\sum y_i$，则 $tx+(1-t)y=\sum(tx_i+(1-t)y_i)$ 属于和；积逐坐标同理。闭包：选 $x_n,y_n\in C$ 趋于闭包两点，凸组合在 $C$，极限在闭包。凸集连通：若被相对开集分离，连接两侧点的连续路径 $t\mapsto tx+(1-t)y$ 把分离拉回 $[0,1]$，与区间连通矛盾。最后若 $C$ 闭凸且含0，$x_n/t_n\to d$ 且 $t_n\to\infty$，充分大时 $x_n/t_n=(1-1/t_n)0+x_n/t_n\in C$，闭性给 $d\in C$。

\Slide{49}{凸性的基本运算及连通性}
本页继续“凸性的基本运算及连通性”。完整定义、条件和论证见本讲原页第 48 页下方。

\Slide{50}{凸性的基本运算及连通性}
本页继续“凸性的基本运算及连通性”。完整定义、条件和论证见本讲原页第 48 页下方。

\Slide{51}{凸性的基本运算及连通性}
本页继续“凸性的基本运算及连通性”。完整定义、条件和论证见本讲原页第 48 页下方。

\Slide{52}{凸性的基本运算及连通性}
本页继续“凸性的基本运算及连通性”。完整定义、条件和论证见本讲原页第 48 页下方。

\Slide{53}{凸性的基本运算及连通性}
本页继续“凸性的基本运算及连通性”。完整定义、条件和论证见本讲原页第 48 页下方。

\Slide{54}{凸性的基本运算及连通性}
本页继续“凸性的基本运算及连通性”。完整定义、条件和论证见本讲原页第 48 页下方。

\Slide{55}{凸包与和的交换及证明纠错}
凸包是所有包含 $S$ 的凸集的交，等价于有限凸组合全体。闭凸包是它的闭包。证明 $\operatorname{co}(A+B)=\operatorname{co}A+\operatorname{co}B$：一边把 $\sum\lambda_i(a_i+b_i)$ 分成两凸组合；另一边若 $a=\sum_i\alpha_i a_i,b=\sum_j\beta_j b_j$，则 $a+b=\sum_{i,j}\alpha_i\beta_j(a_i+b_j)$，积权重非负且总和1。有限多集合用权重 $\prod_k\lambda_{ki_k}$。第58页直接把总权重 $K$ 归一化会把目标点也除以 $K$，因此不成立；第59页积权重论证才正确。

\Slide{56}{凸包与和的交换及证明纠错}
本页继续“凸包与和的交换及证明纠错”。完整定义、条件和论证见本讲原页第 55 页下方。

\Slide{57}{凸包与和的交换及证明纠错}
本页继续“凸包与和的交换及证明纠错”。完整定义、条件和论证见本讲原页第 55 页下方。

\Slide{58}{凸包与和的交换及证明纠错}
本页继续“凸包与和的交换及证明纠错”。完整定义、条件和论证见本讲原页第 55 页下方。

\Slide{59}{凸包与和的交换及证明纠错}
本页继续“凸包与和的交换及证明纠错”。完整定义、条件和论证见本讲原页第 55 页下方。

\Slide{60}{内积、超平面与半空间}
$p\cdot x=\sum_i p_ix_i$ 为内积。固定非零 $p$，$H=\{x:p\cdot x=c\}$ 为超平面，$p$ 为法向量；非零条件避免整个空间或空集。两闭半空间由 $\ge c$ 与 $\le c$ 定义。线性函数连续，弱不等式在极限下保持，故闭；线性组合保持同侧不等式，故凸。

\Slide{61}{内积、超平面与半空间}
本页继续“内积、超平面与半空间”。完整定义、条件和论证见本讲原页第 60 页下方。

\Slide{62}{内积、超平面与半空间}
本页继续“内积、超平面与半空间”。完整定义、条件和论证见本讲原页第 60 页下方。

\Slide{63}{支撑超平面的完整证明}
要证凸集 $K$ 在非内点 $x$ 有过 $x$ 的界超平面。先对非空闭凸 $C$ 与外点 $z$ 求最近点 $a$：最小化距离可限制到一个闭有界球，与 $C$ 的交紧，故最小存在。对任意 $y\in C$，线段 $a+t(y-a)$ 可行，比较平方距离并除以 $t>0$ 再令 $t\downarrow0$，得 $(z-a)\cdot(y-a)\le0$。若 $x\notin C$，取 $p=x-a$，则 $p\cdot(y-x)\le-\norm p^2<0$，过 $x$ 的面已界住 $C$。若 $x\in\partial C$，取外点 $z_n\to x$ 及最近点 $a_n$，有 $\norm{a_n-x}\le2\norm{z_n-x}\to0$。单位法向 $p_n=(z_n-a_n)/\norm{z_n-a_n}$ 抽子列趋于单位 $p$，对每个 $y\in C$ 取极限得 $p\cdot(y-x)\le0$。一般凸 $K$ 令 $C=\overline K$。若内域为空，则 $K$ 的仿射包为真子空间（否则有 $N+1$ 个仿射独立点，其单纯形有内域），取与其方向正交的非零法向即可过任意 $x$ 界住 $K$。若内域非空，则 $\operatorname{int}\overline K=\operatorname{int}K$：闭包内点周围可在 $K$ 选取逼近一组包住该点的小单纯形顶点，凸性使更小邻域在 $K$，故仍可用前述两情形。必要性：若 $x$ 为内点，沿 $p$ 的正负小位移均在 $K$，不可能在面同侧。空 $K$ 时任意过 $x$ 的面都满足。

\Slide{64}{支撑超平面的完整证明}
本页继续“支撑超平面的完整证明”。完整定义、条件和论证见本讲原页第 63 页下方。

\Slide{65}{极锥与坐标序}
$C^\circ=\{p:p\cdot x\le0\ \forall x\in C\}$ 是各个闭半空间的交，故闭凸，且非负倍乘仍满足不等式，所以是锥。正正交象限 $\R_+^N$ 的极锥是 $\R_-^N$：取各坐标单位向量证明必要，逐项相乘求和证明充分。坐标序的 $<$ 只要求至少一坐标严格，$\ll$ 要求全部严格。

\Slide{66}{极锥与坐标序}
本页继续“极锥与坐标序”。完整定义、条件和论证见本讲原页第 65 页下方。

\Slide{67}{三种无穷远方向及量词}
截尾定义先保留 $\norm x\ge M$ 的点，再取所生成最小闭锥，最后对全部 $M$ 取交。序列定义为 $d\in C_\infty\iff\exists x_n\in C,\ t_n\to+\infty:x_n/t_n\to d$。衰退锥为 $\operatorname{rec}C=\{d:\forall x\in C,\forall t\ge0,x+td\in C\}$，要求所有起点都可以一直沿该方向走；只在非空闭凸情形才保证与渐近锥一致。单位方向归一化必须选范数趋于无穷的子列，而不仅是声称任意无界序列范数都趋无穷。截尾最小闭锥不是凸锥，闭包横线不能漏。

\Slide{68}{三种无穷远方向及量词}
本页继续“三种无穷远方向及量词”。完整定义、条件和论证见本讲原页第 67 页下方。

\Slide{69}{三种无穷远方向及量词}
本页继续“三种无穷远方向及量词”。完整定义、条件和论证见本讲原页第 67 页下方。

\Slide{70}{三种无穷远方向及量词}
本页继续“三种无穷远方向及量词”。完整定义、条件和论证见本讲原页第 67 页下方。

\Slide{71}{三种无穷远方向及量词}
本页继续“三种无穷远方向及量词”。完整定义、条件和论证见本讲原页第 67 页下方。

\Slide{72}{三种无穷远方向及量词}
本页继续“三种无穷远方向及量词”。完整定义、条件和论证见本讲原页第 67 页下方。

\Slide{73}{三种无穷远方向及量词}
本页继续“三种无穷远方向及量词”。完整定义、条件和论证见本讲原页第 67 页下方。

\Slide{74}{三种无穷远方向及量词}
本页继续“三种无穷远方向及量词”。完整定义、条件和论证见本讲原页第 67 页下方。

\Slide{75}{三种无穷远方向及量词}
本页继续“三种无穷远方向及量词”。完整定义、条件和论证见本讲原页第 67 页下方。

\Slide{76}{渐近锥闭性及闭包不变}
非空 $C$ 时固定 $x\in C$ 并取 $t_n=n$ 得0属于渐近锥。正倍数通过把 $t_n$ 换为 $t_n/\lambda$ 得到。闭性用对角选取：若 $d_j\to d$，对每个 $j$ 从其见证序列选一点使 $t_j>j$ 且 $\norm{x_j/t_j-d_j}<1/j$；三角不等式给 $x_j/t_j\to d$。闭包不改变锥：用 $x_j'\in C$ 逼近 $x_j\in\overline C$ 至误差 $1/j$，除以 $t_j$ 后误差更趋0。若 $C$ 本身是闭锥，取 $x_n=nd,t_n=n$ 得 $C\subset C_\infty$，反向用每个 $x_n/t_n\in C$ 与闭性；非闭锥的结论为 $C_\infty=\overline C$。

\Slide{77}{渐近锥闭性及闭包不变}
本页继续“渐近锥闭性及闭包不变”。完整定义、条件和论证见本讲原页第 76 页下方。

\Slide{78}{渐近锥闭性及闭包不变}
本页继续“渐近锥闭性及闭包不变”。完整定义、条件和论证见本讲原页第 76 页下方。

\Slide{79}{渐近锥闭性及闭包不变}
本页继续“渐近锥闭性及闭包不变”。完整定义、条件和论证见本讲原页第 76 页下方。

\Slide{80}{渐近锥闭性及闭包不变}
本页继续“渐近锥闭性及闭包不变”。完整定义、条件和论证见本讲原页第 76 页下方。

\Slide{81}{截尾与序列定义等价}
非零 $d=x_n/t_n$ 的极限使 $\norm{x_n}\to\infty$，所以对任意 $M$，最终原点均属截尾集，其缩放极限落入截尾闭锥，故属于交。反向对非零 $d$，每个整数 $M$ 选 $\alpha_M x_M$ 距 $d$ 小于 $1/M$，且 $\norm{x_M}\ge M$。由 $\alpha_M\le(\norm d+1/M)/M\to0$，取 $t_M=1/\alpha_M\to\infty$，得到见证序列。零向量对非空集合单独处理。归一化方向全体再取非负倍数得到相同锥；有界集方向集为空时需单独按约定加入0。

\Slide{82}{截尾与序列定义等价}
本页继续“截尾与序列定义等价”。完整定义、条件和论证见本讲原页第 81 页下方。

\Slide{83}{截尾与序列定义等价}
本页继续“截尾与序列定义等价”。完整定义、条件和论证见本讲原页第 81 页下方。

\Slide{84}{有界判据与平移不变}
非空有界 $C$ 的任意 $x_n/t_n$ 趋于0，故渐近锥只有0。若无界，选 $\norm{x_n}\to\infty$，单位向量抽子列趋于单位 $d$，令 $t_n=\norm{x_n}$ 得非零渐近方向，矛盾。平移 $C+a$ 不改变方向，因为 $(x_n+a)/t_n-x_n/t_n=a/t_n\to0$。若 $E$ 非空，固定 $a\in E$，由 $C+a\subset C+E$ 及单调性得 $C_\infty\subset(C+E)_\infty$。

\Slide{85}{有界判据与平移不变}
本页继续“有界判据与平移不变”。完整定义、条件和论证见本讲原页第 84 页下方。

\Slide{86}{有界判据与平移不变}
本页继续“有界判据与平移不变”。完整定义、条件和论证见本讲原页第 84 页下方。

\Slide{87}{有界判据与平移不变}
本页继续“有界判据与平移不变”。完整定义、条件和论证见本讲原页第 84 页下方。

\Slide{88}{乘积与交集的渐近方向}
若 $(x_{1n},\ldots,x_{mn})/t_n\to(d_1,\ldots,d_m)$，逐坐标就是 $x_{in}/t_n\to d_i$，所以积的渐近锥包含于渐近锥之积。反向一般不能单靠各自见证序列，因为缩放速度未必相同。第89页把总范数截尾当成各分量都截尾的等式错误，例如 $(M,0)$ 总范数大但第二分量不大。交集包含于每个集合，锥的单调性给交集锥包含于各锥交；若右边只有0，交集有界（空交集本身有界，非空时用第84页）。

\Slide{89}{乘积与交集的渐近方向}
本页继续“乘积与交集的渐近方向”。完整定义、条件和论证见本讲原页第 88 页下方。

\Slide{90}{乘积与交集的渐近方向}
本页继续“乘积与交集的渐近方向”。完整定义、条件和论证见本讲原页第 88 页下方。

\Slide{91}{乘积与交集的渐近方向}
本页继续“乘积与交集的渐近方向”。完整定义、条件和论证见本讲原页第 88 页下方。

\Slide{92}{线性像定理再证明}
本段重复第37--41页的重要定理。核心核条件为 $\ker L\cap C_\infty=\{0\}$。若 $Lx_n$ 有界而 $x_n$ 无界，取范数趋无穷子列并单位化，极限是非零渐近方向且被 $L$ 送到0，矛盾。因此像序列收敛时原像可抽收敛子列，闭 $C$ 与连续 $L$ 给像闭。锥像的容易包含是对 $x_n/t_n$ 直接取 $L$；困难包含须用第41页的缩放反证，不可仅引用原集合上的有界性。

\Slide{93}{线性像定理再证明}
本页继续“线性像定理再证明”。完整定义、条件和论证见本讲原页第 92 页下方。

\Slide{94}{线性像定理再证明}
本页继续“线性像定理再证明”。完整定义、条件和论证见本讲原页第 92 页下方。

\Slide{95}{线性像定理再证明}
本页继续“线性像定理再证明”。完整定义、条件和论证见本讲原页第 92 页下方。

\Slide{96}{线性像定理再证明}
本页继续“线性像定理再证明”。完整定义、条件和论证见本讲原页第 92 页下方。

\Slide{97}{闭和定理的两个证明}
积集加求和映射的证明已在第42页完整给出。两集合的直接证明：若 $z_n=a_n+b_n\to z$ 而 $a_n$ 无界，选 $r_n=\norm{a_n}\to\infty$ 及 $a_n/r_n\to d$，$\norm d=1$。由 $z_n/r_n\to0$ 得 $b_n/r_n\to-d$，于是 $d\in A_\infty,-d\in B_\infty$，违反正半独立。故 $a_n$ 有界，抽子列趋于 $a\in A$，再由 $b_n=z_n-a_n\to z-a\in B$，得 $z\in A+B$。有限多集合用积方法最直接，避免在归纳时漏检中间和的渐近锥条件。

\Slide{98}{闭和定理的两个证明}
本页继续“闭和定理的两个证明”。完整定义、条件和论证见本讲原页第 97 页下方。

\Slide{99}{闭和定理的两个证明}
本页继续“闭和定理的两个证明”。完整定义、条件和论证见本讲原页第 97 页下方。

\Slide{100}{闭和定理的两个证明}
本页继续“闭和定理的两个证明”。完整定义、条件和论证见本讲原页第 97 页下方。

\Slide{101}{闭和定理的两个证明}
本页继续“闭和定理的两个证明”。完整定义、条件和论证见本讲原页第 97 页下方。

\Slide{102}{闭和定理的两个证明}
本页继续“闭和定理的两个证明”。完整定义、条件和论证见本讲原页第 97 页下方。

\Slide{103}{闭和定理的两个证明}
本页继续“闭和定理的两个证明”。完整定义、条件和论证见本讲原页第 97 页下方。

\Slide{104}{渐近正则性与凸渐近锥}
正则性把“存在缩放速度”增强为“每一种趋无穷速度都可匹配”。设 $x_j/s_j\to d$，固定 $a\in C$，给任意 $t_n\to\infty$ 选 $j(n)\to\infty$ 且 $s_{j(n)}\ge t_n$。由凸性 $z_n=a+(t_n/s_{j(n)})(x_{j(n)}-a)\in C$，且 $z_n/t_n=a/t_n+(x_{j(n)}-a)/s_{j(n)}\to d$。因此任意非空凸集渐近正则。对两方向使用共同 $t_n=n$ 的见证点，再作相同权重凸组合，就证明渐近锥凸；其闭锥性质已由第76页给出。

\Slide{105}{渐近正则性与凸渐近锥}
本页继续“渐近正则性与凸渐近锥”。完整定义、条件和论证见本讲原页第 104 页下方。

\Slide{106}{渐近正则性与凸渐近锥}
本页继续“渐近正则性与凸渐近锥”。完整定义、条件和论证见本讲原页第 104 页下方。

\Slide{107}{渐近正则性与凸渐近锥}
本页继续“渐近正则性与凸渐近锥”。完整定义、条件和论证见本讲原页第 104 页下方。

\Slide{108}{渐近正则性与凸渐近锥}
本页继续“渐近正则性与凸渐近锥”。完整定义、条件和论证见本讲原页第 104 页下方。

\Slide{109}{衰退锥相等与非闭反例}
任意非空集合有 $\operatorname{rec}C\subset C_\infty$：固定 $a\in C$，用 $x_n=a+nd$。反向需要闭凸：若 $x_n/t_n\to d$，对任意 $a\in C$ 和固定 $t>0$，充分大时 $a+(t/t_n)(x_n-a)\in C$，取极限得 $a+td\in C$。所以 $\operatorname{rec}C=C_\infty$。第112页反例 $C=([0,1)\times[0,\infty))\cup\{(1,0)\}$ 中，衰退方向必须让 $(1,0)$ 及任意 $(0,y)$ 都留在集合，迫使两个方向分量均0，故衰退锥是 $\{0\}$，不是空集；渐近锥是 $\{(0,v):v\ge0\}$，不是整条竖直线。闭凸时 $C+C_\infty=C$ 由衰退定义及0在锥中立即得到。

\Slide{110}{衰退锥相等与非闭反例}
本页继续“衰退锥相等与非闭反例”。完整定义、条件和论证见本讲原页第 109 页下方。

\Slide{111}{衰退锥相等与非闭反例}
本页继续“衰退锥相等与非闭反例”。完整定义、条件和论证见本讲原页第 109 页下方。

\Slide{112}{衰退锥相等与非闭反例}
本页继续“衰退锥相等与非闭反例”。完整定义、条件和论证见本讲原页第 109 页下方。

\Slide{113}{衰退锥相等与非闭反例}
本页继续“衰退锥相等与非闭反例”。完整定义、条件和论证见本讲原页第 109 页下方。

\Slide{114}{最大可含锥与基点无关}
闭凸 $S$、$a\in S$ 下，令 $D_a=\bigcap_{\lambda>0}\lambda(S-a)$。$d\in D_a$ 等价于任意 $t>0$ 有 $a+td\in S$，因此它是从 $a$ 出发可完整放入 $S$ 的最大锥。若 $b\in S$，取 $T\to\infty$，由凸性 $(1-t/T)b+(t/T)(a+Td)\in S$，极限为 $b+td$，闭性给后者在 $S$，所以基点可以任意换，$D_a=\operatorname{rec}S=S_\infty$。当0在 $S$ 中，取 $a=0,t=1$ 就得到渐近锥包含于原集。本段所有闭性使用均在最后取极限处，不是多余假设。

\Slide{115}{最大可含锥与基点无关}
本页继续“最大可含锥与基点无关”。完整定义、条件和论证见本讲原页第 114 页下方。

\Slide{116}{最大可含锥与基点无关}
本页继续“最大可含锥与基点无关”。完整定义、条件和论证见本讲原页第 114 页下方。

\Slide{117}{最大可含锥与基点无关}
本页继续“最大可含锥与基点无关”。完整定义、条件和论证见本讲原页第 114 页下方。

\Slide{118}{最大可含锥与基点无关}
本页继续“最大可含锥与基点无关”。完整定义、条件和论证见本讲原页第 114 页下方。
